\section{The GHZ state}\label{sec:ghz}
%% ===================================================================

The Greenberger--Horne--Zeilinger (GHZ) state is the prototypical
non-injective, renormalization-fixed-point MPS.

\begin{definition}[GHZ tensor]\label{def:ghz_tensor}
    \lean{MPSTensor.ghzTensor}
    \leanok
    \uses{def:mps_tensor}
    Set $d = D = 2$. The GHZ tensor is the family
    $\{A^0, A^1\} \subset \MN{2}$ defined by
    \[
        A^0 = \ketbra{0}{0} = \diag(1,0),
        \qquad
        A^1 = \ketbra{1}{1} = \diag(0,1).
    \]
    For a system of $N$ sites the resulting MPV is the GHZ state
    \[
        \ket{\mathrm{GHZ}_N}
        = \ket{0}^{\otimes N} + \ket{1}^{\otimes N}.
    \]
\end{definition}

\begin{definition}[Two-site computational vector]\label{def:two_site_ket}
    \lean{MPSTensor.twoSiteKet}
    \leanok
    For $a,b\in\{0,1\}$, the vector
    \[
        \ket{ab}\in(\{0,1\}^{2}\to\C)
    \]
    is the indicator of the configuration
    \[
        (\sigma_1,\sigma_2)=(a,b).
    \]
\end{definition}

\begin{definition}[Constant product vector]\label{def:constant_ket}
    \lean{MPSTensor.constantKet}
    \leanok
    For $a\in\{0,1\}$ and $N\ge0$,
    \[
        \ket{a}^{\otimes N}(\sigma_1,\ldots,\sigma_N)
        =
        \begin{cases}
            1, & (\sigma_1,\ldots,\sigma_N)=(a,\ldots,a), \\
            0, & \text{otherwise}.
        \end{cases}
    \]
\end{definition}

\begin{theorem}[GHZ transfer map is idempotent]\label{thm:ghz_transfer_idempotent}
    \lean{MPSTensor.ghz_transferMap_idempotent}
    \leanok
    \uses{def:ghz_tensor, def:transfer_map}
    The transfer map of the GHZ tensor satisfies $\E_A^2 = \E_A$.
\end{theorem}

\begin{proof}
    \leanok
    The transfer map acts by
    $\E_A(X)_{ij} = \delta_{ij} X_{ij}$
    (extraction of the diagonal).
    This operation is idempotent.
\end{proof}

\begin{theorem}[GHZ tensor is an RFP]\label{thm:ghz_is_rfp}
    \lean{MPSTensor.ghz_isTransferIdempotent}
    \leanok
    \uses{def:ghz_tensor, def:mps_transfer_idempotence,
        thm:ghz_transfer_idempotent}
    The GHZ tensor is a renormalization fixed point.
\end{theorem}

\begin{proof}
    \leanok
    \uses{thm:ghz_transfer_idempotent}
    The claim follows from Theorem~\ref{thm:ghz_transfer_idempotent}.
\end{proof}

\begin{theorem}[GHZ tensor is not injective]\label{thm:ghz_not_injective}
    \lean{MPSTensor.ghz_not_isInjective}
    \leanok
    \uses{def:ghz_tensor, def:injective}
    The GHZ tensor is not injective: $\spn_{\C}\{A^0, A^1\}$ is the
    two-dimensional space of diagonal matrices in~$\MN{2}$, which does
    not equal~$\MN{2}$.
\end{theorem}

\begin{proof}
    \leanok
    Every element of $\spn_{\C}\{A^0, A^1\}$ has zero off-diagonal
    entries, so the matrix with all entries equal to~$1$ is not in the
    span.
\end{proof}

\begin{theorem}[$\Z_2$ on-site symmetry of the GHZ tensor]\label{thm:ghz_z2_symmetric}
    \lean{MPSTensor.ghz_isOnSiteSymmetric_Z2}
    \leanok
    \uses{def:ghz_tensor, def:on_site_symmetric}
    Let $\Z_2 = \{1, \sigma_x\}$ act on $\C^2$ by the Pauli-$X$
    matrix~$\sigma_x$. The GHZ tensor is on-site symmetric under this
    action: for each $g \in \Z_2$, the twisted tensor $\widetilde A_g$
    defines the same MPV family as $A$, i.e.\ $\mc V(A)=\mc V(\widetilde
        A_g)$.
\end{theorem}

\begin{proof}
    \leanok
    The identity element is trivial. For the generator
    $g = \sigma_x$, the twisted tensor satisfies
    $\widetilde A_g^i = A^{1-i}$, i.e.\ the two matrices are swapped.
    Conjugation by $\sigma_x$ implements the same swap, giving
    gauge equivalence and hence the same MPV family.
\end{proof}

\begin{definition}[Virtual $\Z_2$ representation of the GHZ tensor]\label{def:ghz_virtual_z2_action}
    \lean{MPSTensor.ghzVirtualZ2Action}
    \leanok
    \uses{def:ghz_tensor}
    The virtual $\Z_2$ representation on the bond space is the homomorphism
    $\Z_2 \to \MN{2}$ sending the identity to $I$ and the generator to
    $\sigma_z = \diag(1, -1)$.
\end{definition}

\begin{theorem}[Virtual $\Z_2$ symmetry of the GHZ tensor]\label{thm:ghz_virtual_z2}
    \lean{MPSTensor.ghz_virtual_Z2_symmetric}
    \leanok
    \uses{def:ghz_tensor, def:ghz_virtual_z2_action}
    The matrix $\sigma_z$ commutes with every GHZ matrix $A^i$.  A non-scalar
    matrix commuting with all $A^i$ is the hallmark of a non-injective
    ($\Z_2$-injective) tensor: the on-site symmetry of
    Theorem~\ref{thm:ghz_z2_symmetric} swaps the two one-dimensional injective
    blocks $A^0 \leftrightarrow A^1$, while $\sigma_z$ acts on each block by a
    phase.
\end{theorem}

\begin{proof}
    \leanok
    Both $A^0$ and $A^1$ are diagonal, so each commutes with the diagonal
    matrix $\sigma_z$.
\end{proof}

\begin{theorem}[GHZ order parameter is a transfer-map fixed point]\label{thm:ghz_order_param}
    \lean{MPSTensor.ghz_orderParam_fixed}
    \leanok
    \uses{def:ghz_tensor, def:transfer_map}
    The order-parameter operator $\sigma_z$ is a fixed point of the GHZ transfer
    map, $\E_A(\sigma_z) = \sigma_z$.
\end{theorem}

\begin{proof}
    \leanok
    Since $\E_A(X) = \sum_i A^i X (A^i)^{\dagger}$ and both $A^i$ are diagonal,
    the claim is the entrywise computation $\E_A(\sigma_z) = \sigma_z$.
\end{proof}

\begin{remark}
    A primitive (injective) tensor has a one-dimensional transfer-map
    fixed-point space, spanned by a positive matrix that equals the identity only
    in a normalized gauge.  That the GHZ transfer map fixes the additional
    non-scalar operator $\sigma_z$ is the transfer-matrix signature of the extra
    non-decaying mode behind its long-range ferromagnetic correlations.
\end{remark}

\begin{theorem}[GHZ zero correlation length]\label{thm:ghz_zcl}
    \lean{MPSTensor.ghz_isZCL}
    \leanok
    \uses{def:is_zcl, thm:ghz_is_rfp, thm:zcl_iff_idempotent_transfer}
    The GHZ tensor has zero correlation length: its idempotent transfer map
    (Theorem~\ref{thm:ghz_is_rfp}) gives correlations independent of the
    separation.  This contrasts with the AKLT state
    (Theorem~\ref{thm:aklt_correlation_length}), whose subleading eigenvalue
    $-\tfrac13$ produces a finite correlation length $1/\log 3$.
\end{theorem}

\begin{proof}
    \leanok
    \uses{thm:zcl_iff_idempotent_transfer, thm:ghz_is_rfp}
    The claim follows from the idempotence of the GHZ transfer map and the equivalence
    of zero correlation length with transfer-map idempotence
    (Theorem~\ref{thm:zcl_iff_idempotent_transfer}).
\end{proof}

%% ===================================================================
